\section{Exploratory Data Analysis}
\label{sec:eda}

Before building models, we characterized the data through a suite of
exploratory visualizations. Five findings shaped the feature engineering
decisions of Section~\ref{sec:methods}.

\textbf{Finding 1: Ridership follows strong daily and weekly cycles.}
Figure~\ref{fig:hour} shows mean passenger count by hour of day, separated
by weekday and weekend. Weekdays exhibit pronounced morning (7--9~AM) and
afternoon (4--7~PM) peaks typical of commuter transit, while weekends show
a broader single midday peak centered around 1--3~PM. This justifies
including \texttt{hour}, \texttt{dayofweek}, and \texttt{is\_rush\_hour}
as calendar-derived temporal features.

\begin{figure}[t]
  \centering
  \includegraphics[width=\columnwidth]{03_ridership_by_hour.png}
  \caption{Mean passenger count by hour of day, separated by weekday and
    weekend. Weekdays show classic bimodal rush-hour structure; weekends
    show a unimodal midday peak.}
  \label{fig:hour}
\end{figure}

\textbf{Finding 2: Seasonal and COVID effects are visible in the longitudinal
data.} Figure~\ref{fig:seasonal} shows monthly mean ridership across the
2019--2022 study period. A sharp drop in mid-2020 reflects the COVID-19
lockdown period in Zurich, followed by a gradual recovery through 2021 and
2022. This motivated the inclusion of COVID intensity features derived from
the Oxford COVID-19 Government Response Tracker (retained in the baseline
feature set but not a focus of this paper).

\begin{figure}[t]
  \centering
  \includegraphics[width=\columnwidth]{04_seasonal_trends.png}
  \caption{Monthly mean passenger count across the study period. The
    pronounced mid-2020 drop reflects the COVID-19 lockdown; the
    2021--2022 gradual recovery is clearly visible.}
  \label{fig:seasonal}
\end{figure}

\textbf{Finding 3: Routes exhibit distinct demand profiles.} The 10 ZTBus
routes have visibly different passenger count distributions: some routes are
consistently busier than others, and some show more temporal variance. We
encode \texttt{busRoute} as a one-hot categorical feature with 9 dummies
(drop-first to avoid collinearity).

\textbf{Finding 4: Feature correlations identify redundancy.}
Figure~\ref{fig:corr} shows the Pearson correlation matrix for the nine
numeric raw and engineered features. Vehicle speed and power demand are
strongly correlated (as expected, since the vehicle is electric and
acceleration demand drives the battery), while ambient temperature has
near-zero correlation with passenger count ($r = -0.011$). Based on this
observation, we exclude raw temperature from the final feature set.

\begin{figure}[t]
  \centering
  \includegraphics[width=\columnwidth]{07_correlation_matrix.png}
  \caption{Pearson correlation matrix of numeric features. Speed and power
    demand are strongly correlated (expected for an electric vehicle);
    temperature has near-zero correlation with passenger count and is
    excluded from the feature set.}
  \label{fig:corr}
\end{figure}

\textbf{Finding 5: Stop names concentrate mass in a small set.} The 20
most-visited stops account for the majority of stop events across the
dataset. We encode \texttt{stopName} by retaining the top-20 categories
and bucketing all other stops as \texttt{\_\_other\_\_}, avoiding a 147-way
categorical explosion.

These observations translate into the six feature groups described in
Section~\ref{sec:methods}: Temporal, Spatial, Operational, Sensor, Status,
and Categorical. The full feature specification is deferred to
Section~\ref{sec:methods} for continuity.
